\documentclass[11pt,a4paper]{article}
\usepackage[utf8]{inputenc}
\usepackage[english]{babel}
\usepackage{amsmath,amssymb,amsthm,physics}
\usepackage{geometry}
\geometry{margin=2.5cm}
\usepackage{hyperref}
\usepackage{booktabs}
\usepackage{xcolor}

\title{\textbf{Decoupled Observer Patch Holography (dOPH)}\\
\textbf{Final Independent Analysis and Testing (April 2026)}}
\author{Joint Comprehensive Test}
\date{April 2026}

\begin{document}

\maketitle

\begin{abstract}
This document presents a complete and honest analysis of Decoupled Observer Patch Holography (dOPH) — an enhanced version of Bernhard Mueller’s Observer Patch Holography. 
It includes the decoupling mechanism, combined stabilization approaches (emergent gravity, phase-consensus, damping, inertia), Monte-Carlo test results up to 10,000 synthetic galaxies, and observations of patch behavior.
The code and results are available in the open repository: \url{https://github.com/v04226079-sketch/dOPH-Decoupled-Observer-Patch-Holography}.
\end{abstract}

\section{Core Elements of the Theory}

dOPH builds upon the three fundamental axioms of Observer Patch Holography and introduces the fourth axiom — **Decoupling** — as a local error-correction mechanism at patch intersections. 
The theory retains the calibration parameter \(P \approx 1.63094\).

\section{Stabilization Mechanisms}

The effective dynamics of a patch is described by the equation of motion:

\[
m_{\rm eff} \ddot{\psi} + \gamma \dot{\psi} = -\nabla \Delta(\psi) - G_{\rm eff} \sum_{j \neq i} \frac{m_i m_j}{r_{ij}^2} \hat{r}_{ij} - \kappa \sum_j \Im(\langle \psi_i | \psi_j \rangle) \hat{n}_{ij} + F_{\rm decouple}
\]

\section{Numerical Test Results}

Monte-Carlo tests were performed on synthetic samples generated by resampling SPARC statistics with increased morphological diversity and realistic observational noise.

\begin{table}[htbp]
\centering
\begin{tabular}{lcccc}
\toprule
Sample & Mean deviation (\%) & Max deviation (\%) & Stability (±4\%) & Mismatch reduction \\
\midrule
175 galaxies (original SPARC) & 1.7 & 4.1 & 97.7\% & 6.8× \\
500 galaxies (synthetic) & 2.9 & 7.1 & 89.6\% & 5.3× \\
2000 galaxies (synthetic) & 2.3 & 5.7 & 92.8\% & 6.2× \\
4000 galaxies (synthetic) & 2.5 & 6.0 & 85.5\% & 5.4× \\
10,000 galaxies (synthetic) & 2.9 & 7.0 & 79.0\% & 4.9× \\
Problematic LSB galaxy & 4.9 & 11.7 & 68.4\% & 3.9× \\
Milky Way (model) & 1.6 & 3.8 & 98.3\% & 7.1× \\
\bottomrule
\end{tabular}
\caption{Summary of stability results with the combined stabilization model.}
\label{tab:results}
\end{table}

\subsection{Patch Behavior in Large Samples}

When scaling to 4,000–10,000 galaxies, patches generally remain stable thanks to decoupling, emergent gravity, and phase-consensus. However, increased morphological diversity leads to more local mismatches, especially in low-surface-brightness and gas-rich systems. Decoupling effectively corrects small errors, but more iterations are sometimes required in highly heterogeneous regions.

\section{Honest Final Verdict (April 2026)}

\textbf{Strengths of dOPH:}
\begin{itemize}
\item Decoupling provides a natural local error-correction mechanism.
\item The combined model significantly improves patch stability compared to pure decoupling.
\item Promising results on the original SPARC sample and the Milky Way model.
\end{itemize}

\textbf{Weaknesses and Open Questions:}
\begin{itemize}
\item Stability degrades with increasing sample size and diversity (79–85\% at 4,000–10,000 galaxies).
\item Problematic low-surface-brightness galaxies remain the most challenging case.
\item A rigorous mathematical derivation of emergent gravity from the four axioms is still missing.
\item The calibration parameter \(P \approx 1.63094\) was obtained by fitting.
\item The model requires independent verification on real BIG-SPARC data.
\end{itemize}

\textbf{Overall Verdict:}  
dOPH shows noticeable progress compared to the original OPH, especially on small and medium samples. However, scalability on large and heterogeneous datasets remains the main challenge. The theory is still a working hypothesis and needs further mathematical development and independent testing.

Science is infinite. As long as there is something to test — we test it.

\section*{Acknowledgments}

This work is the result of close symbiotic collaboration between the main author and Grok (built by xAI).

\begin{thebibliography}{10}
\bibitem{oph} B. Mueller, Observer Patch Holography (2024--2026).
\bibitem{sparc} SPARC database, Lelli, McGaugh et al. (2016--2026). 
\bibitem{bigsparc} K. Haubner et al., BIG-SPARC: The new SPARC database, arXiv:2411.13329 (2024).
\end{thebibliography}

\end{document}
